Independently trained self-play agents often adopt arbitrary conventions and then fail to coordinate with each other. We reproduce this in a small CPU study with exactly evaluated payoffs on three symmetric games (a lever game, a safe-option matrix game, a signalling game) and compare self-play (SP), other-play (OP) and a reimplemented fictitious-co-play-style population method on cross-play between independently trained agents (20 agents $\times$ 5 seeds). SP reaches self-play payoff of at least 0.97 but cross-play of only 0.08--0.27. OP raises cross-play on the safe and signalling games, but there it does so by picking the safe action and by abandoning communication; in the lever game its cross-play ($0.15$ vs $0.08$) is not distinguishable from SP at five seeds ($p=0.053$) because most OP agents still collapse onto an arbitrary lever, and this is insensitive to training length, and OP's cross-play tends to fall with larger initial logit scale (not significant at five seeds) as fewer agents start inside the basin of the distinct lever. The population method lies between SP and OP in the mean on the safe and signalling games (OP vs population on safe: $p=0.08$), improves with population size, and gives no gain in the lever game. Results are small-scale and tabular.
